\documentclass[aspectratio=169,10pt,spanish]{beamer}

\input{../../../_preambulo_beamer.tex}
\input{../../../_comandos_beamer.tex}

\selectlanguage{spanish}

\title{MA1001B - Modelación Estadística para la Toma de Decisiones}
\subtitle{Lección 30: Tamaño de muestra para margen de error y presupuesto}
\author{}
\institute{Tecnol\'ogico de Monterrey}
\date{}

\begin{document}

\begin{frame}[plain]
  \vspace{1.2cm}
  {\Large\bfseries MA1001B - Modelación Estadística para la Toma de Decisiones\par}
  \vspace{0.7cm}
  {\large Lección 30: Tamaño de muestra para margen de error y presupuesto\par}
  \vspace{0.7cm}
  {\color{TecRojo}\rule{\textwidth}{1.2pt}}
  \vspace{0.8cm}

  Facultad MA1001B\\[0.3cm]
  Periodo\\[0.3cm]
  Tecnol\'ogico de Monterrey
\end{frame}

\begin{frame}{La pregunta de planificación}
  \begin{alertblock}{Pregunta guía}
    ¿Qué tan grande debe ser $n$ para que un intervalo al 95\% alcance un
    margen de error elegido, y puede el presupuesto financiar ese $n$?
  \end{alertblock}

  \vspace{0.6em}
  Al final de esta lección, usted puede:
  \begin{itemize}
    \item calcular $n=(z^{*}s/E)^{2}$ para una media y redondear hacia arriba;
    \item calcular $n=z^{*2}p(1-p)/E^{2}$ para una proporción;
    \item explicar por qué $p=0.5$ es conservadora;
    \item comparar el $n$ requerido con el $n$ que un presupuesto puede comprar.
  \end{itemize}

  \vfill
  \scriptsize Todos los valores de planificación usados hoy son sintéticos. No se usan datos reales de empresa.
\end{frame}

\begin{frame}{Elegir $E$ antes de recolectar datos}
  \small
  \begin{center}
    \begin{tabular}{lr}
      \toprule
      \textbf{Plan de tiempo medio de atención} & \textbf{Valor} \\
      \midrule
      $z^{*}$ de libro & 1.96 \\
      $s$ de planificación & 10 minutos \\
      Margen objetivo $E$ & 2 minutos \\
      \bottomrule
    \end{tabular}
  \end{center}

  \vspace{0.5em}
  \begin{exampleblock}{Fórmula de tamaño de muestra para la media}
    \[
      n=\left(\frac{z^{*}s}{E}\right)^{2}=\left(\frac{1.96\times 10}{2}\right)^{2}=96.040000.
    \]
    Redondear hacia arriba a la siguiente observación entera: $n=97$.
  \end{exampleblock}
\end{frame}

\begin{frame}[fragile]{Paso 1: Tamaño de muestra para una media}
  \begin{columns}[T]
    \begin{column}{0.42\textwidth}
      \small
      \begin{lstlisting}
raw = (z_star * s / E) ** 2
n = int(np.ceil(raw))
      \end{lstlisting}

      \begin{block}{Salida verificada}
        $n$ crudo $=96.040000$\\
        $n$ redondeado $=97$
      \end{block}
    \end{column}
    \begin{column}{0.58\textwidth}
      \centering
      \includegraphics[width=\textwidth]{../figuras/sample_size_01_mean_margin.png}
      \vspace{0.2em}
      \scriptsize $n$ requerido versus $E$ del Paso 1
    \end{column}
  \end{columns}
\end{frame}

\begin{frame}{$p=0.5$ conservadora para una proporción}
  \small
  Margen objetivo $E=0.03$ para una tasa de primer contacto:
  \[
    n=\frac{z^{*2}p(1-p)}{E^{2}}.
  \]
  En $p=0.5$, $p(1-p)=0.25$ se maximiza:
  \[
    n=\frac{1.96^{2}\times 0.25}{0.03^{2}}=1067.111111 \;\to\; 1068.
  \]
  Usar la suposición $p=0.18$ de la Lección 29 produce solo $n=631$, que es
  437 observaciones demasiado optimista si $p$ está más cerca de $0.5$.
\end{frame}

\begin{frame}[fragile]{Paso 2: Tamaño de muestra conservador para una proporción}
  \begin{columns}[T]
    \begin{column}{0.40\textwidth}
      \small
      \begin{lstlisting}
raw = (
    z_star**2 * p * (1-p)
    / E**2
)
      \end{lstlisting}

      \begin{block}{Salida verificada}
        $p=0.5$ da $n=1068$\\
        $p=0.18$ da $n=631$\\
        Extra $=437$
      \end{block}
    \end{column}
    \begin{column}{0.60\textwidth}
      \centering
      \includegraphics[width=\textwidth]{../figuras/sample_size_02_proportion_conservative.png}
      \vspace{0.2em}
      \scriptsize $n$ versus $p$ de planificación del Paso 2
    \end{column}
  \end{columns}
\end{frame}

\begin{frame}{Paso 3: El $E=0.01$ deseado puede exceder el presupuesto}
  \begin{columns}[T]
    \begin{column}{0.42\textwidth}
      \small
      $n$ conservador para $E=0.01$:
      \[
        n=9604.
      \]
      Un presupuesto de $\$20{,}000$ a $\$12$ cada una
      financia solo $1666$ observaciones.

      \vspace{0.4em}
      Costo requerido: $\$115{,}248$.

      \vspace{0.5em}
      \textbf{Límite:} reducir $E$ infla $n$ con $1/E^{2}$. Un margen
      deseado no es un estudio financiable.
    \end{column}
    \begin{column}{0.58\textwidth}
      \centering
      \includegraphics[width=\textwidth]{../figuras/sample_size_03_budget_vs_precision.png}
      \vspace{0.2em}
      \scriptsize $n$ requerido versus asequible del Paso 3
    \end{column}
  \end{columns}
\end{frame}

\begin{frame}{Verificación de conceptos}
  \begin{enumerate}
    \item Calcule $(1.96\times 10/2)^{2}$ y redondee hacia arriba.
    \item ¿Por qué $p=0.5$ es conservadora para una proporción?
    \item ¿Por qué $n$ explota cuando $E$ se reduce de $0.03$ a $0.01$?
    \item ¿Puede un presupuesto de $\$20{,}000$ a $\$12$ por observación
      financiar $n=9604$?
  \end{enumerate}

  \vspace{0.6em}
  \begin{block}{Conexión con la Lección 31}
    La sesión puede continuar directamente con un intervalo ji-cuadrada
    para una varianza poblacional.
  \end{block}

  \vfill
  \scriptsize Fuente: OpenStax, \textit{Introductory Business Statistics 2e},
  Capítulo 8, Sección 8.4. CC BY-NC-SA.
\end{frame}

\appendix



\end{document}
